%%%%%%%%%%%%%%%%%%%%%%%%%%%%%%%%%%%%%%%%%%%%%%%%%%%%%%%%%%%%%%%%%%%%%%%
% Universidade do Estado de Santa Catarina           
% Graduação em Engenharia de Software                  
% (c) 2013 Eduardo Henrique Couto (ehcouto@hotmail.com)
% (c) 2015 Adaptado por Arthur Garcia Bartsch (arthurbartsch@gmail.com)
% (c) 2019 Adaptado por Paolo Moser (paolo.moser@gmail.com), Tiago Luiz Schmitz e Fernando dos Santos 

%%%%%%%%%%%%%%%%%%%%%%%%%%%%%%%%%%%%%%%%%%%%%%%%%%%%%%%%%%%%%%%%%%%%%%%
\documentclass{./Data/udesc} % Classe da udesc  

 %%%%%%%%%%%%%%%%%%%%%%%%%%%%%%%%%%%%%%%%%%%%%%%%%%%%%%%%%%%%%%%%%%%%%%
\usepackage{./Data/udescsty} % Estilo da UDESC
\usepackage{./Data/udescdata} % Dados do trabalho
\usepackage{./Data/udescpretextual} % Elementos pré-textuais do trabalho
\usepackage{./Data/udescpackages} % Outros pacotes do trabalho

%%%%%%%%%%%%%%%%%%%%%%%%%%%%%%%%%%%%%%%%%%%%%%%%%%%%%%%%%%%%%%%%%%%%%%%
% ORIENTAÇÕES PARA AUTORES (ALUNOS)
% Para definir os dados do texto, consulte o arquivo udescdata.sty
% Para definir os elementos pré-textuais, consulte o arquivo udescpretextual.sty
% Para definir os pacotes, ver alguns comandos especiais (entre outros), consulte o arquivo udescpackages.sty
% Procure não alterar os demais arquivos da pasta "data"
%%%%%%%%%%%%%%%%%%%%%%%%%%%%%%%%%%%%%%%%%%%%%%%%%%%%%%%%%%%%%%%%%%%%%%%

\begin{document}
%--------------------------------------------------------
% Elementos pré-textuais

\pretextual % Ativa todos os elementos pré-textuais

%--------------------------------------------------------
% Elementos textuais

% Inclua os capítulos aqui
\include{./tex/Introducao}
\include{./tex/Conclusao}

%--------------------------------------------------------
% Elementos pós-textuais

% Use as referências geradas pelo Google Scholar - certifique-se de que todos os campos estão corretos!
\renewcommand{\refname}{REFERÊNCIAS}
\bibliographystyle{./Data/abntex2-alf}
\bibliography{./Ref/Referencias}

%%% Inclua apendices ou anexos
%\apendice
%\include{}
%\anexo
%\include{}
\end{document}